\documentclass[10pt]{extarticle}
\usepackage[letterpaper, margin=1in]{geometry}
\usepackage{amsmath}
\usepackage{amssymb}
\usepackage{enumitem}
\usepackage{fancyhdr}
\usepackage{lastpage}
\usepackage{tcolorbox}

\allowdisplaybreaks
\setlength{\parindent}{0pt}

\title{\centering\Large\textsc{Applications of Linear Algebra: Problem Set 2}}
\author{\large \textsc{Rento Saijo}}
\date{March 26, 2026}

\pagestyle{fancy}
\fancyhf{}
\fancyfoot[C]{\textsc{Page \thepage\ of \pageref{LastPage}}}
\renewcommand{\headrulewidth}{0pt}

\newtcolorbox{problem}[1][]{
    colframe=darkgray!100,
    arc=0.1mm,
    boxrule=2pt,
    boxsep=1mm,
    title={\textsc{#1}}
}

\begin{document}
\maketitle
\thispagestyle{fancy}
\vspace{-1em}\noindent\rule{\linewidth}{0.6pt}\vspace{0.5em}

\begin{problem}[Problem 1]
Recall that a matrix is called unitriangular if it is triangular with all $1$'s on the diagonal. Show that the product of two lower unitriangular matrices is again lower unitriangular. It might help to first show that the product of two lower triangular matrices is again lower triangular, and then consider the entries on the diagonal.
\end{problem}

Let
\[
A=(a_{ij}),\qquad B=(b_{ij}),\qquad C=AB=(c_{ij}),
\]
and assume both $A$ and $B$ are lower unitriangular. The entry formula for a matrix product is
\begin{equation}
c_{ij}=\sum_{k=1}^n a_{ik}b_{kj}.
\label{eq:product-entry}
\end{equation}
So the whole problem is really about understanding which terms in \eqref{eq:product-entry} can actually survive. First, we want to show that $C$ is lower triangular. Take any entry above the diagonal, so $j>i$. In the sum \eqref{eq:product-entry}, two things can happen:
\begin{itemize}[nosep]
\item if $k>i$, then $a_{ik}=0$ because $A$ is lower triangular;
\item if $k\le i<j$, then $b_{kj}=0$ because $B$ is lower triangular.
\end{itemize}
So every term $a_{ik}b_{kj}$ is zero, which means
\[
c_{ij}=0 \qquad \text{whenever } j>i.
\]
That proves $C$ is lower triangular. Now, for the diagonal. If $i=j$, then \eqref{eq:product-entry} becomes
\[
c_{ii}=\sum_{k=1}^n a_{ik}b_{ki}.
\]
Again, almost all of those terms vanish:
\begin{itemize}[nosep]
\item if $k>i$, then $a_{ik}=0$;
\item if $k<i$, then $b_{ki}=0$.
\end{itemize}
So only the $k=i$ term remains, and therefore
\begin{equation}
c_{ii}=a_{ii}b_{ii}=1\cdot 1=1.
\label{eq:unit-diagonal}
\end{equation}

So $C$ is lower triangular, and by \eqref{eq:unit-diagonal}, every diagonal entry is $1$. Therefore, the product $AB$ is again lower unitriangular.

\newpage

\begin{problem}[Problem 2]
Let $t_1,t_2,\dots$ be distinct real numbers. The $n\times n$ Vandermonde matrix, $V(t_1,t_2,\dots,t_n)$ is the matrix of the form
\[
V(t_1,t_2,\dots,t_n)=
\begin{pmatrix}
1 & 1 & 1 & \cdots & 1\\
t_1 & t_2 & t_3 & \cdots & t_n\\
t_1^2 & t_2^2 & t_3^2 & \cdots & t_n^2\\
\vdots & \vdots & \vdots & \ddots & \vdots\\
t_1^{n-1} & t_2^{n-1} & t_3^{n-1} & \cdots & t_n^{n-1}
\end{pmatrix}.
\]
These matrices show up in lots of applications such as the design matrix of polynomial regression, the defining matrix of discrete Fourier transforms, and the normalized volumes of certain polytopes. I don't think you should necessarily know what any of that means, but I hope you recognize that those are wildly different fields where Vandermonde matrices show up. Find the LU factorization of the $2\times 2$, $3\times 3$, and $4\times 4$ Vandermonde matrices. What is the pattern you see? Test your pattern with the $5\times 5$ Vandermonde matrix.
\end{problem}

For the $2\times 2$ case, start with
\[
V(t_1,t_2)=
\begin{pmatrix}
1 & 1\\
t_1 & t_2
\end{pmatrix}.
\]
\[
R_2\leftarrow R_2-t_1R_1.
\]
The corresponding elementary matrix is
\[
E_1=
\begin{pmatrix}
1 & 0\\
-t_1 & 1
\end{pmatrix}.
\]
So
\[
E_1V(t_1,t_2)=
\begin{pmatrix}
1 & 0\\
-t_1 & 1
\end{pmatrix}
\begin{pmatrix}
1 & 1\\
t_1 & t_2
\end{pmatrix}
=
\begin{pmatrix}
1 & 1\\
0 & t_2-t_1
\end{pmatrix}
=U.
\]
The inverse elementary matrix is
\[
L_1=E_1^{-1}=
\begin{pmatrix}
1 & 0\\
t_1 & 1
\end{pmatrix},
\]
and therefore,
\[
V(t_1,t_2)=L_1U=
\begin{pmatrix}
1 & 0\\
t_1 & 1
\end{pmatrix}
\begin{pmatrix}
1 & 1\\
0 & t_2-t_1
\end{pmatrix}.
\]

For the $3\times 3$ case, let
\[
\!A_0=V(t_1,t_2,t_3)=
\begin{pmatrix}
1 & 1 & 1\\
t_1 & t_2 & t_3\\
t_1^2 & t_2^2 & t_3^2
\end{pmatrix}.
\]
First, use
\[
R_2\leftarrow R_2-t_1R_1.
\]
The corresponding elementary matrix is
\[
E_1=
\begin{pmatrix}
1 & 0 & 0\\
-t_1 & 1 & 0\\
0 & 0 & 1
\end{pmatrix}.
\]
So
\[
A_1=E_1A_0=
\begin{pmatrix}
1 & 1 & 1\\
0 & t_2-t_1 & t_3-t_1\\
t_1^2 & t_2^2 & t_3^2
\end{pmatrix}.
\]

Next, use
\[
R_3\leftarrow R_3-t_1^2R_1.
\]
The corresponding elementary matrix is
\[
E_2=
\begin{pmatrix}
1 & 0 & 0\\
0 & 1 & 0\\
-t_1^2 & 0 & 1
\end{pmatrix}.
\]
So
\[
A_2=E_2A_1=E_2E_1A_0=
\begin{pmatrix}
1 & 1 & 1\\
0 & t_2-t_1 & t_3-t_1\\
0 & t_2^2-t_1^2 & t_3^2-t_1^2
\end{pmatrix}.
\]

Now, use
\[
R_3\leftarrow R_3-(t_1+t_2)R_2.
\]
The corresponding elementary matrix is
\[
E_3=
\begin{pmatrix}
1 & 0 & 0\\
0 & 1 & 0\\
0 & -(t_1+t_2) & 1
\end{pmatrix}.
\]
Then,
\[
A_3=E_3A_2=E_3E_2E_1A_0
=
\begin{pmatrix}
1 & 1 & 1\\
0 & t_2-t_1 & t_3-t_1\\
0 & 0 & (t_3-t_1)(t_3-t_2)
\end{pmatrix}
=U.
\]

Now, we invert the elementary matrices:
\[
L_1=E_1^{-1}=
\begin{pmatrix}
1 & 0 & 0\\
t_1 & 1 & 0\\
0 & 0 & 1
\end{pmatrix},
\qquad
L_2=E_2^{-1}=
\begin{pmatrix}
1 & 0 & 0\\
0 & 1 & 0\\
t_1^2 & 0 & 1
\end{pmatrix},
\qquad
L_3=E_3^{-1}=
\begin{pmatrix}
1 & 0 & 0\\
0 & 1 & 0\\
0 & t_1+t_2 & 1
\end{pmatrix}.
\]
Their product is
\[
L=L_1L_2L_3=E_1^{-1}E_2^{-1}E_3^{-1}=
\begin{pmatrix}
1 & 0 & 0\\
t_1 & 1 & 0\\
t_1^2 & t_1+t_2 & 1
\end{pmatrix},
\]
so
\[
V(t_1,t_2,t_3)=LU=
\begin{pmatrix}
1 & 0 & 0\\
t_1 & 1 & 0\\
t_1^2 & t_1+t_2 & 1
\end{pmatrix}
\begin{pmatrix}
1 & 1 & 1\\
0 & t_2-t_1 & t_3-t_1\\
0 & 0 & (t_3-t_1)(t_3-t_2)
\end{pmatrix}.
\]

For the $4\times 4$ case, let
\[
A_0=V(t_1,t_2,t_3,t_4)=
\begin{pmatrix}
1 & 1 & 1 & 1\\
t_1 & t_2 & t_3 & t_4\\
t_1^2 & t_2^2 & t_3^2 & t_4^2\\
t_1^3 & t_2^3 & t_3^3 & t_4^3
\end{pmatrix}.
\]
Similar to the above cases, let us perform the operations in order.

\[
R_2\leftarrow R_2-t_1R_1,
\qquad
E_1=
\begin{pmatrix}
1 & 0 & 0 & 0\\
-t_1 & 1 & 0 & 0\\
0 & 0 & 1 & 0\\
0 & 0 & 0 & 1
\end{pmatrix},
\]
\[
A_1=E_1A_0=
\begin{pmatrix}
1 & 1 & 1 & 1\\
0 & t_2-t_1 & t_3-t_1 & t_4-t_1\\
t_1^2 & t_2^2 & t_3^2 & t_4^2\\
t_1^3 & t_2^3 & t_3^3 & t_4^3
\end{pmatrix}.
\]

\[
R_3\leftarrow R_3-t_1^2R_1,
\qquad
E_2=
\begin{pmatrix}
1 & 0 & 0 & 0\\
0 & 1 & 0 & 0\\
-t_1^2 & 0 & 1 & 0\\
0 & 0 & 0 & 1
\end{pmatrix},
\]
\[
A_2=E_2A_1=
\begin{pmatrix}
1 & 1 & 1 & 1\\
0 & t_2-t_1 & t_3-t_1 & t_4-t_1\\
0 & t_2^2-t_1^2 & t_3^2-t_1^2 & t_4^2-t_1^2\\
t_1^3 & t_2^3 & t_3^3 & t_4^3
\end{pmatrix}.
\]

\[
R_4\leftarrow R_4-t_1^3R_1,
\qquad
E_3=
\begin{pmatrix}
1 & 0 & 0 & 0\\
0 & 1 & 0 & 0\\
0 & 0 & 1 & 0\\
-t_1^3 & 0 & 0 & 1
\end{pmatrix},
\]
\[
A_3=E_3A_2=
\begin{pmatrix}
1 & 1 & 1 & 1\\
0 & t_2-t_1 & t_3-t_1 & t_4-t_1\\
0 & t_2^2-t_1^2 & t_3^2-t_1^2 & t_4^2-t_1^2\\
0 & t_2^3-t_1^3 & t_3^3-t_1^3 & t_4^3-t_1^3
\end{pmatrix}.
\]

\[
R_3\leftarrow R_3-(t_1+t_2)R_2,
\qquad
E_4=
\begin{pmatrix}
1 & 0 & 0 & 0\\
0 & 1 & 0 & 0\\
0 & -(t_1+t_2) & 1 & 0\\
0 & 0 & 0 & 1
\end{pmatrix},
\]
\[
A_4=E_4A_3=
\begin{pmatrix}
1 & 1 & 1 & 1\\
0 & t_2-t_1 & t_3-t_1 & t_4-t_1\\
0 & 0 & (t_3-t_1)(t_3-t_2) & (t_4-t_1)(t_4-t_2)\\
0 & t_2^3-t_1^3 & t_3^3-t_1^3 & t_4^3-t_1^3
\end{pmatrix}.
\]

\[
R_4\leftarrow R_4-(t_1^2+t_1t_2+t_2^2)R_2,
\qquad
E_5=
\begin{pmatrix}
1 & 0 & 0 & 0\\
0 & 1 & 0 & 0\\
0 & 0 & 1 & 0\\
0 & -(t_1^2+t_1t_2+t_2^2) & 0 & 1
\end{pmatrix},
\]
\[
A_5=E_5A_4=
\begin{pmatrix}
1 & 1 & 1 & 1\\
0 & t_2-t_1 & t_3-t_1 & t_4-t_1\\
0 & 0 & (t_3-t_1)(t_3-t_2) & (t_4-t_1)(t_4-t_2)\\
0 & 0 & (t_3-t_1)(t_3-t_2)(t_1+t_2+t_3) & (t_4-t_1)(t_4-t_2)(t_1+t_2+t_4)
\end{pmatrix}.
\]

\[
R_4\leftarrow R_4-(t_1+t_2+t_3)R_3,
\qquad
E_6=
\begin{pmatrix}
1 & 0 & 0 & 0\\
0 & 1 & 0 & 0\\
0 & 0 & 1 & 0\\
0 & 0 & -(t_1+t_2+t_3) & 1
\end{pmatrix},
\]
\[
A_6=E_6A_5=
\begin{pmatrix}
1 & 1 & 1 & 1\\
0 & t_2-t_1 & t_3-t_1 & t_4-t_1\\
0 & 0 & (t_3-t_1)(t_3-t_2) & (t_4-t_1)(t_4-t_2)\\
0 & 0 & 0 & (t_4-t_1)(t_4-t_2)(t_4-t_3)
\end{pmatrix}
=U.
\]

\begingroup
\small
\[
L_1=E_1^{-1}=
\begin{pmatrix}
1 & 0 & 0 & 0\\
t_1 & 1 & 0 & 0\\
0 & 0 & 1 & 0\\
0 & 0 & 0 & 1
\end{pmatrix},
\quad
L_2=E_2^{-1}=
\begin{pmatrix}
1 & 0 & 0 & 0\\
0 & 1 & 0 & 0\\
t_1^2 & 0 & 1 & 0\\
0 & 0 & 0 & 1
\end{pmatrix},
\quad
L_3=E_3^{-1}=
\begin{pmatrix}
1 & 0 & 0 & 0\\
0 & 1 & 0 & 0\\
0 & 0 & 1 & 0\\
t_1^3 & 0 & 0 & 1
\end{pmatrix},
\]
\[
L_4=E_4^{-1}=
\begin{pmatrix}
1 & 0 & 0 & 0\\
0 & 1 & 0 & 0\\
0 & t_1+t_2 & 1 & 0\\
0 & 0 & 0 & 1
\end{pmatrix},
\]
\[
L_5=E_5^{-1}=
\begin{pmatrix}
1 & 0 & 0 & 0\\
0 & 1 & 0 & 0\\
0 & 0 & 1 & 0\\
0 & t_1^2+t_1t_2+t_2^2 & 0 & 1
\end{pmatrix},
\quad
L_6=E_6^{-1}=
\begin{pmatrix}
1 & 0 & 0 & 0\\
0 & 1 & 0 & 0\\
0 & 0 & 1 & 0\\
0 & 0 & t_1+t_2+t_3 & 1
\end{pmatrix}.
\]
\endgroup
Their product is
\begingroup
\[
L=
\begin{pmatrix}
1 & 0 & 0 & 0\\
t_1 & 1 & 0 & 0\\
t_1^2 & t_1+t_2 & 1 & 0\\
t_1^3 & t_1^2+t_1t_2+t_2^2 & t_1+t_2+t_3 & 1
\end{pmatrix},
\]
\[
V(t_1,t_2,t_3,t_4)=LU.
\]
\endgroup

At this point, the pattern is pretty clear. The upper-triangular entries look like
\begin{equation}
u_{ij}=\prod_{m=1}^{i-1}(t_j-t_m)\qquad \text{for } i\le j,
\label{eq:vander-u}
\end{equation}
with $u_{ij}=0$ for $i>j$. For the lower-triangular factor, the cleanest way to describe the pattern is column by column. In column 1, the entries are just powers of $t_1$:
\[
1,\ t_1,\ t_1^2,\ t_1^3,\ t_1^4,\dots
\]
In column 2, each entry is the sum of all monomials of the needed degree in $t_1$ and $t_2$:
\[
1,\ t_1+t_2,\ t_1^2+t_1t_2+t_2^2,\ t_1^3+t_1^2t_2+t_1t_2^2+t_2^3,\dots
\]
In column 3, we do the same thing, but now with $t_1,t_2,t_3$:
\[
1,\ t_1+t_2+t_3,\ t_1^2+t_2^2+t_3^2+t_1t_2+t_1t_3+t_2t_3,\dots
\]
So, in general, the entry $l_{ij}$ is found by taking all monomials of total degree $i-j$ in the variables $t_1,\dots,t_j$ and adding them together. For example,
\[
l_{52}=t_1^3+t_1^2t_2+t_1t_2^2+t_2^3,
\qquad
l_{53}=t_1^2+t_2^2+t_3^2+t_1t_2+t_1t_3+t_2t_3.
\]
That tells us what we should expect in the $5\times5$ case before we even factorize. In particular, we expect the last row of $L$ to be
\[
\bigl(t_1^4,\,
t_1^3+t_1^2t_2+t_1t_2^2+t_2^3,\,
t_1^2+t_2^2+t_3^2+t_1t_2+t_1t_3+t_2t_3,\,
t_1+t_2+t_3+t_4,\,
1\bigr),
\]
and we expect the last diagonal entry of $U$ to be $(t_5-t_1)(t_5-t_2)(t_5-t_3)(t_5-t_4)$. Now we check that prediction directly in the $5\times5$ case. Let
\begingroup
\small
\[
A_0=V(t_1,t_2,t_3,t_4,t_5)=
\begin{pmatrix}
1 & 1 & 1 & 1 & 1\\
t_1 & t_2 & t_3 & t_4 & t_5\\
t_1^2 & t_2^2 & t_3^2 & t_4^2 & t_5^2\\
t_1^3 & t_2^3 & t_3^3 & t_4^3 & t_5^3\\
t_1^4 & t_2^4 & t_3^4 & t_4^4 & t_5^4
\end{pmatrix}.
\]

To keep the later entries readable, let
\[
q_3(t)=t_1+t_2+t,\qquad
r_3(t)=t_1^2+t_2^2+t^2+t_1t_2+t_1t+t_2t,\qquad
q_4(t)=t_1+t_2+t_3+t.
\]

\[
R_2\leftarrow R_2-t_1R_1,
\qquad
A_1=
\begin{pmatrix}
1 & 1 & 1 & 1 & 1\\
0 & t_2-t_1 & t_3-t_1 & t_4-t_1 & t_5-t_1\\
t_1^2 & t_2^2 & t_3^2 & t_4^2 & t_5^2\\
t_1^3 & t_2^3 & t_3^3 & t_4^3 & t_5^3\\
t_1^4 & t_2^4 & t_3^4 & t_4^4 & t_5^4
\end{pmatrix},
\]

\[
R_3\leftarrow R_3-t_1^2R_1,
\qquad
A_2=
\begin{pmatrix}
1 & 1 & 1 & 1 & 1\\
0 & t_2-t_1 & t_3-t_1 & t_4-t_1 & t_5-t_1\\
0 & t_2^2-t_1^2 & t_3^2-t_1^2 & t_4^2-t_1^2 & t_5^2-t_1^2\\
t_1^3 & t_2^3 & t_3^3 & t_4^3 & t_5^3\\
t_1^4 & t_2^4 & t_3^4 & t_4^4 & t_5^4
\end{pmatrix},
\]

\[
R_4\leftarrow R_4-t_1^3R_1,
\qquad
A_3=
\begin{pmatrix}
1 & 1 & 1 & 1 & 1\\
0 & t_2-t_1 & t_3-t_1 & t_4-t_1 & t_5-t_1\\
0 & t_2^2-t_1^2 & t_3^2-t_1^2 & t_4^2-t_1^2 & t_5^2-t_1^2\\
0 & t_2^3-t_1^3 & t_3^3-t_1^3 & t_4^3-t_1^3 & t_5^3-t_1^3\\
t_1^4 & t_2^4 & t_3^4 & t_4^4 & t_5^4
\end{pmatrix},
\]

\[
R_5\leftarrow R_5-t_1^4R_1,
\qquad
A_4=
\begin{pmatrix}
1 & 1 & 1 & 1 & 1\\
0 & t_2-t_1 & t_3-t_1 & t_4-t_1 & t_5-t_1\\
0 & t_2^2-t_1^2 & t_3^2-t_1^2 & t_4^2-t_1^2 & t_5^2-t_1^2\\
0 & t_2^3-t_1^3 & t_3^3-t_1^3 & t_4^3-t_1^3 & t_5^3-t_1^3\\
0 & t_2^4-t_1^4 & t_3^4-t_1^4 & t_4^4-t_1^4 & t_5^4-t_1^4
\end{pmatrix},
\]

\[
R_3\leftarrow R_3-(t_1+t_2)R_2,
\qquad
A_5=
\begin{pmatrix}
1 & 1 & 1 & 1 & 1\\
0 & t_2-t_1 & t_3-t_1 & t_4-t_1 & t_5-t_1\\
0 & 0 & (t_3-t_1)(t_3-t_2) & (t_4-t_1)(t_4-t_2) & (t_5-t_1)(t_5-t_2)\\
0 & t_2^3-t_1^3 & t_3^3-t_1^3 & t_4^3-t_1^3 & t_5^3-t_1^3\\
0 & t_2^4-t_1^4 & t_3^4-t_1^4 & t_4^4-t_1^4 & t_5^4-t_1^4
\end{pmatrix},
\]

\[
R_4\leftarrow R_4-(t_1^2+t_1t_2+t_2^2)R_2.
\]
\[
A_6=
\begin{pmatrix}
1 & 1 & 1 & 1 & 1\\
0 & t_2-t_1 & t_3-t_1 & t_4-t_1 & t_5-t_1\\
0 & 0 & (t_3-t_1)(t_3-t_2) & (t_4-t_1)(t_4-t_2) & (t_5-t_1)(t_5-t_2)\\
0 & 0 & (t_3-t_1)(t_3-t_2)q_3(t_3) & (t_4-t_1)(t_4-t_2)q_3(t_4) & (t_5-t_1)(t_5-t_2)q_3(t_5)\\
0 & t_2^4-t_1^4 & t_3^4-t_1^4 & t_4^4-t_1^4 & t_5^4-t_1^4
\end{pmatrix},
\]

\[
R_5\leftarrow R_5-(t_1^3+t_1^2t_2+t_1t_2^2+t_2^3)R_2.
\]
\[
A_7=
\begin{pmatrix}
1 & 1 & 1 & 1 & 1\\
0 & t_2-t_1 & t_3-t_1 & t_4-t_1 & t_5-t_1\\
0 & 0 & (t_3-t_1)(t_3-t_2) & (t_4-t_1)(t_4-t_2) & (t_5-t_1)(t_5-t_2)\\
0 & 0 & (t_3-t_1)(t_3-t_2)q_3(t_3) & (t_4-t_1)(t_4-t_2)q_3(t_4) & (t_5-t_1)(t_5-t_2)q_3(t_5)\\
0 & 0 & (t_3-t_1)(t_3-t_2)r_3(t_3) & (t_4-t_1)(t_4-t_2)r_3(t_4) & (t_5-t_1)(t_5-t_2)r_3(t_5)
\end{pmatrix},
\]

\[
R_4\leftarrow R_4-(t_1+t_2+t_3)R_3.
\]
\[
A_8=
\begin{pmatrix}
1 & 1 & 1 & 1 & 1\\
0 & t_2-t_1 & t_3-t_1 & t_4-t_1 & t_5-t_1\\
0 & 0 & (t_3-t_1)(t_3-t_2) & (t_4-t_1)(t_4-t_2) & (t_5-t_1)(t_5-t_2)\\
0 & 0 & 0 & (t_4-t_1)(t_4-t_2)(t_4-t_3) & (t_5-t_1)(t_5-t_2)(t_5-t_3)\\
0 & 0 & (t_3-t_1)(t_3-t_2)r_3(t_3) & (t_4-t_1)(t_4-t_2)r_3(t_4) & (t_5-t_1)(t_5-t_2)r_3(t_5)
\end{pmatrix},
\]

\[
R_5\leftarrow R_5-(t_1^2+t_2^2+t_3^2+t_1t_2+t_1t_3+t_2t_3)R_3.
\]
\[
A_9=
\begin{pmatrix}
1 & 1 & 1 & 1 & 1\\
0 & t_2-t_1 & t_3-t_1 & t_4-t_1 & t_5-t_1\\
0 & 0 & (t_3-t_1)(t_3-t_2) & (t_4-t_1)(t_4-t_2) & (t_5-t_1)(t_5-t_2)\\
0 & 0 & 0 & (t_4-t_1)(t_4-t_2)(t_4-t_3) & (t_5-t_1)(t_5-t_2)(t_5-t_3)\\
0 & 0 & 0 & (t_4-t_1)(t_4-t_2)(t_4-t_3)q_4(t_4) & (t_5-t_1)(t_5-t_2)(t_5-t_3)q_4(t_5)
\end{pmatrix},
\]

\[
R_5\leftarrow R_5-(t_1+t_2+t_3+t_4)R_4.
\]
\[
A_{10}=
\begin{pmatrix}
1 & 1 & 1 & 1 & 1\\
0 & t_2-t_1 & t_3-t_1 & t_4-t_1 & t_5-t_1\\
0 & 0 & (t_3-t_1)(t_3-t_2) & (t_4-t_1)(t_4-t_2) & (t_5-t_1)(t_5-t_2)\\
0 & 0 & 0 & (t_4-t_1)(t_4-t_2)(t_4-t_3) & (t_5-t_1)(t_5-t_2)(t_5-t_3)\\
0 & 0 & 0 & 0 & (t_5-t_1)(t_5-t_2)(t_5-t_3)(t_5-t_4)
\end{pmatrix}
=U.
\]

\[
L=
\begin{pmatrix}
1 & 0 & 0 & 0 & 0\\
t_1 & 1 & 0 & 0 & 0\\
t_1^2 & t_1+t_2 & 1 & 0 & 0\\
t_1^3 & t_1^2+t_1t_2+t_2^2 & t_1+t_2+t_3 & 1 & 0\\
t_1^4 & t_1^3+t_1^2t_2+t_1t_2^2+t_2^3 & t_1^2+t_2^2+t_3^2+t_1t_2+t_1t_3+t_2t_3 & t_1+t_2+t_3+t_4 & 1
\end{pmatrix},
\]
\[
V(t_1,t_2,t_3,t_4,t_5)=LU.
\]
\endgroup

This matches the pattern in \eqref{eq:vander-u}, and it also confirms the lower-triangular pattern we predicted before carrying out the $5\times5$ factorization.

\newpage

\begin{problem}[Problem 3]
\begin{enumerate}[label=(\alph*)]
\item Explain why the following are all legitimate permuted LU factorizations of the same matrix. List the elementary row operations that are being used in each case.
\begin{align*}
\begin{pmatrix}
0 & 1 & 0\\
1 & 0 & 0\\
0 & 0 & 1
\end{pmatrix}
\begin{pmatrix}
0 & 1 & 3\\
2 & -1 & 1\\
2 & -2 & 0
\end{pmatrix}
&=
\begin{pmatrix}
1 & 0 & 0\\
0 & 1 & 0\\
1 & -1 & 1
\end{pmatrix}
\begin{pmatrix}
2 & -1 & 1\\
0 & 1 & 3\\
0 & 0 & 2
\end{pmatrix}\\
\begin{pmatrix}
0 & 0 & 1\\
0 & 1 & 0\\
1 & 0 & 0
\end{pmatrix}
\begin{pmatrix}
0 & 1 & 3\\
2 & -1 & 1\\
2 & -2 & 0
\end{pmatrix}
&=
\begin{pmatrix}
1 & 0 & 0\\
1 & 1 & 0\\
0 & 1 & 1
\end{pmatrix}
\begin{pmatrix}
2 & -2 & 0\\
0 & 1 & 1\\
0 & 0 & 2
\end{pmatrix}\\
\begin{pmatrix}
0 & 0 & 1\\
1 & 0 & 0\\
0 & 1 & 0
\end{pmatrix}
\begin{pmatrix}
0 & 1 & 3\\
2 & -1 & 1\\
2 & -2 & 0
\end{pmatrix}
&=
\begin{pmatrix}
1 & 0 & 0\\
0 & 1 & 0\\
1 & 1 & 1
\end{pmatrix}
\begin{pmatrix}
2 & -2 & 0\\
0 & 1 & 3\\
0 & 0 & -2
\end{pmatrix}
\end{align*}
\item Use each of the factorizations to solve the linear system
\[
\begin{pmatrix}
0 & 1 & 3\\
2 & -1 & 1\\
2 & -2 & 0
\end{pmatrix}
\begin{pmatrix}
x\\
y\\
z
\end{pmatrix}
=
\begin{pmatrix}
-5\\
-1\\
0
\end{pmatrix}.
\]
Do you get the same result in each case? Why or why not?
\end{enumerate}
\end{problem}

Let
\[
A=
\begin{pmatrix}
0 & 1 & 3\\
2 & -1 & 1\\
2 & -2 & 0
\end{pmatrix},
\qquad
\vec{b}=
\begin{pmatrix}
-5\\
-1\\
0
\end{pmatrix}.
\]

\begin{enumerate}[label=(\alph*)]
\item Each identity has the form $PA=LU$, so the permutation matrix $P$ first reorders the rows of $A$, the elementary row operations turn $PA$ into the upper-triangular matrix $U$, and the displayed lower-triangular matrix $L$ is the product of the inverse elementary matrices. So they are all legitimate because they are all different permuted elimination procedures applied to the same matrix. Let us take a look at each of them. For the first factorization, the permutation swaps rows $1$ and $2$, so
\[
P_1A=
\begin{pmatrix}
2 & -1 & 1\\
0 & 1 & 3\\
2 & -2 & 0
\end{pmatrix}.
\]
From there the row operations are
\[
R_3\leftarrow R_3-R_1,\qquad R_3\leftarrow R_3+R_2.
\]
After the first one,
\[
\begin{pmatrix}
2 & -1 & 1\\
0 & 1 & 3\\
0 & -1 & -1
\end{pmatrix},
\]
and after the second one,
\[
\begin{pmatrix}
2 & -1 & 1\\
0 & 1 & 3\\
0 & 0 & 2
\end{pmatrix}
=U_1.
\]
So this really is an elimination procedure that turns $P_1A$ into the displayed upper-triangular matrix. The inverse elementary matrices then multiply together to give
\[
L_1=
\begin{pmatrix}
1 & 0 & 0\\
0 & 1 & 0\\
1 & -1 & 1
\end{pmatrix}.
\]

For the second factorization, the permutation swaps rows $1$ and $3$, so
\[
P_2A=
\begin{pmatrix}
2 & -2 & 0\\
2 & -1 & 1\\
0 & 1 & 3
\end{pmatrix}.
\]
The row operations are
\[
R_2\leftarrow R_2-R_1,\qquad R_3\leftarrow R_3-R_2.
\]
After the first one,
\[
\begin{pmatrix}
2 & -2 & 0\\
0 & 1 & 1\\
0 & 1 & 3
\end{pmatrix},
\]
and after the second one,
\[
\begin{pmatrix}
2 & -2 & 0\\
0 & 1 & 1\\
0 & 0 & 2
\end{pmatrix}
=U_2.
\]
Again, this is exactly the displayed upper-triangular matrix, and the inverse elementary matrices multiply to
\[
L_2=
\begin{pmatrix}
1 & 0 & 0\\
1 & 1 & 0\\
0 & 1 & 1
\end{pmatrix}.
\]

For the third factorization, the permutation reorders the rows as $(3,1,2)$, so
\[
P_3A=
\begin{pmatrix}
2 & -2 & 0\\
0 & 1 & 3\\
2 & -1 & 1
\end{pmatrix}.
\]
The row operations are
\[
R_3\leftarrow R_3-R_1,\qquad R_3\leftarrow R_3-R_2.
\]
After the first one,
\[
\begin{pmatrix}
2 & -2 & 0\\
0 & 1 & 3\\
0 & 1 & 1
\end{pmatrix},
\]
and after the second one,
\[
\begin{pmatrix}
2 & -2 & 0\\
0 & 1 & 3\\
0 & 0 & -2
\end{pmatrix}
=U_3.
\]
So once more the displayed $U$ is exactly what the row operations produce, and the inverse elementary matrices multiply to
\[
L_3=
\begin{pmatrix}
1 & 0 & 0\\
0 & 1 & 0\\
1 & 1 & 1
\end{pmatrix}.
\]

\item We will solve each one using the forward and backward substitution steps: first solve $L\vec{c}=P\vec{b}$, then solve $U\vec{x}=\vec{c}$. For the first factorization,
\[
P_1\vec{b}=
\begin{pmatrix}
-1\\
-5\\
0
\end{pmatrix}.
\]
Solving
\[
\begin{pmatrix}
1 & 0 & 0\\
0 & 1 & 0\\
1 & -1 & 1
\end{pmatrix}
\vec{c}^{(1)}
=
\begin{pmatrix}
-1\\
-5\\
0
\end{pmatrix}
\]
gives
\[
\vec{c}^{(1)}=
\begin{pmatrix}
-1\\
-5\\
-4
\end{pmatrix}.
\]
Then
\[
\begin{pmatrix}
2 & -1 & 1\\
0 & 1 & 3\\
0 & 0 & 2
\end{pmatrix}
\begin{pmatrix}
x\\
y\\
z
\end{pmatrix}
=
\begin{pmatrix}
-1\\
-5\\
-4
\end{pmatrix}
\]
gives $z=-2$, then $y=1$, and then $x=1$. For the second factorization,
\[
P_2\vec{b}=
\begin{pmatrix}
0\\
-1\\
-5
\end{pmatrix}.
\]
Solving
\[
\begin{pmatrix}
1 & 0 & 0\\
1 & 1 & 0\\
0 & 1 & 1
\end{pmatrix}
\vec{c}^{(2)}
=
\begin{pmatrix}
0\\
-1\\
-5
\end{pmatrix}
\]
gives
\[
\vec{c}^{(2)}=
\begin{pmatrix}
0\\
-1\\
-4
\end{pmatrix}.
\]
Then
\[
\begin{pmatrix}
2 & -2 & 0\\
0 & 1 & 1\\
0 & 0 & 2
\end{pmatrix}
\begin{pmatrix}
x\\
y\\
z
\end{pmatrix}
=
\begin{pmatrix}
0\\
-1\\
-4
\end{pmatrix}
\]
again gives $z=-2$, $y=1$, and $x=1$. For the third factorization,
\[
P_3\vec{b}=
\begin{pmatrix}
0\\
-5\\
-1
\end{pmatrix}.
\]
Solving
\[
\begin{pmatrix}
1 & 0 & 0\\
0 & 1 & 0\\
1 & 1 & 1
\end{pmatrix}
\vec{c}^{(3)}
=
\begin{pmatrix}
0\\
-5\\
-1
\end{pmatrix}
\]
gives
\[
\vec{c}^{(3)}=
\begin{pmatrix}
0\\
-5\\
4
\end{pmatrix}.
\]
Then
\[
\begin{pmatrix}
2 & -2 & 0\\
0 & 1 & 3\\
0 & 0 & -2
\end{pmatrix}
\begin{pmatrix}
x\\
y\\
z
\end{pmatrix}
=
\begin{pmatrix}
0\\
-5\\
4
\end{pmatrix}
\]
gives $z=-2$, $y=1$, and $x=1$ one more time. So yes, all three factorizations give the same solution,
\[
\boxed{
\begin{pmatrix}
x\\
y\\
z
\end{pmatrix}
=
\begin{pmatrix}
1\\
1\\
-2
\end{pmatrix}}.
\]
They have to agree, because each system $P_iA\vec{x}=P_i\vec{b}$ is obtained from $A\vec{x}=\vec{b}$ by multiplying both sides by an invertible permutation matrix. That changes the order of the equations, but not the actual solution set.
\end{enumerate}

\newpage

\begin{problem}[Problem 4]
Although the matrix
\[
A=
\begin{pmatrix}
1 & 2 & 2\\
2 & 4 & -1\\
1 & -1 & 3
\end{pmatrix}
\]
is symmetric, show that it cannot be factored as $A=LDL^T$.
\end{problem}

For now, the general fact we want to record is this: any matrix of the form $LDL^T$ must be symmetric, because $D$ is diagonal and therefore $D^T=D$. So
\begin{equation}
(LDL^T)^T=(L^T)^TD^TL^T=LDL^T.
\label{eq:ldlt-symmetric}
\end{equation}
Equation \eqref{eq:ldlt-symmetric} says every $LDL^T$ matrix equals its own transpose, so symmetry is a necessary condition for an $LDL^T$ factorization.

\newpage

\begin{problem}[Problem 5]
Consider the linear system
\[
x-5y-z=1,\qquad \frac{1}{6}x-\frac{5}{6}y+z=0,\qquad 2x-y=3.
\]
\begin{enumerate}[label=(\alph*)]
\item Find the exact solution to the linear system.
\item Solve the system with Gaussian elimination with $4$-digit rounding.
\item Solve the system using Partial Pivoting and $4$-digit rounding.
\item Compare your answers.
\end{enumerate}
\end{problem}

\begin{enumerate}[label=(\alph*)]
\item For the exact solution, let us freely use fractions. The corresponding augmented matrix to the linear system is
\[
\left[
\begin{array}{ccc|c}
1 & -5 & -1 & 1\\
\frac{1}{6} & -\frac{5}{6} & 1 & 0\\
2 & -1 & 0 & 3
\end{array}
\right].
\]
If we use the rows in this order, the second pivot becomes awkward right away, so we will swap the first and third rows first:
\[
R_1\leftrightarrow R_3,
\]
which gives
\[
\left[
\begin{array}{ccc|c}
2 & -1 & 0 & 3\\
\frac{1}{6} & -\frac{5}{6} & 1 & 0\\
1 & -5 & -1 & 1
\end{array}
\right].
\]
Let us get it down to the RREF:
\[
R_2\leftarrow R_2-\frac{1}{12}R_1,\qquad
R_3\leftarrow R_3-\frac{1}{2}R_1,
\]
\[
\left[
\begin{array}{ccc|c}
2 & -1 & 0 & 3\\
0 & -\frac{3}{4} & 1 & -\frac{1}{4}\\
0 & -\frac{9}{2} & -1 & -\frac{1}{2}
\end{array}
\right].
\]
\[
R_3\leftarrow R_3-6R_2,
\]
\[
\left[
\begin{array}{ccc|c}
2 & -1 & 0 & 3\\
0 & -\frac{3}{4} & 1 & -\frac{1}{4}\\
0 & 0 & -7 & 1
\end{array}
\right].
\]
\[
R_3\leftarrow -\frac{1}{7}R_3,
\]
\[
\left[
\begin{array}{ccc|c}
2 & -1 & 0 & 3\\
0 & -\frac{3}{4} & 1 & -\frac{1}{4}\\
0 & 0 & 1 & -\frac{1}{7}
\end{array}
\right].
\]
\[
R_2\leftarrow R_2-R_3,
\]
\[
\left[
\begin{array}{ccc|c}
2 & -1 & 0 & 3\\
0 & -\frac{3}{4} & 0 & -\frac{3}{28}\\
0 & 0 & 1 & -\frac{1}{7}
\end{array}
\right].
\]
\[
R_2\leftarrow -\frac{4}{3}R_2,
\]
\[
\left[
\begin{array}{ccc|c}
2 & -1 & 0 & 3\\
0 & 1 & 0 & \frac{1}{7}\\
0 & 0 & 1 & -\frac{1}{7}
\end{array}
\right].
\]
\[
R_1\leftarrow R_1+R_2,
\]
\[
\left[
\begin{array}{ccc|c}
2 & 0 & 0 & \frac{22}{7}\\
0 & 1 & 0 & \frac{1}{7}\\
0 & 0 & 1 & -\frac{1}{7}
\end{array}
\right].
\]
\[
R_1\leftarrow \frac{1}{2}R_1,
\]
\[
\left[
\begin{array}{ccc|c}
1 & 0 & 0 & \frac{11}{7}\\
0 & 1 & 0 & \frac{1}{7}\\
0 & 0 & 1 & -\frac{1}{7}
\end{array}
\right].
\]
So the exact solution is
\begin{equation}
\boxed{
\begin{pmatrix}
x\\
y\\
z
\end{pmatrix}
=
\begin{pmatrix}
\frac{11}{7}\\[0.3em]
\frac{1}{7}\\[0.3em]
-\frac{1}{7}
\end{pmatrix}}.
\label{eq:exact-solution}
\end{equation}

\item For the rounded computation, we will use $4$ significant digits at each arithmetic step. Here we do \emph{not} pivot, so we simply start with the first row as the first pivot row. The augmented matrix is
\[
\left[
\begin{array}{ccc|c}
1.000 & -5.000 & -1.000 & 1.000\\
0.1667 & -0.8333 & 1.000 & 0\\
2.000 & -1.000 & 0 & 3.000
\end{array}
\right].
\]

\[
R_2\leftarrow R_2-0.1667R_1,\qquad R_3\leftarrow R_3-2.000R_1,
\]
which gives
\[
\left[
\begin{array}{ccc|c}
1.000 & -5.000 & -1.000 & 1.000\\
0 & 0.0002 & 1.167 & -0.1667\\
0 & 9.000 & 2.000 & 1.000
\end{array}
\right].
\]

This is the part where the trouble starts. In exact arithmetic that $0.0002$ would really be $0$, but with rounding it turns into a tiny pivot. So the next row operation is
\[
R_3\leftarrow R_3-(4.500\times 10^4)R_2
\]
and this gives
\begin{equation}
\left[
\begin{array}{ccc|c}
1.000 & -5.000 & -1.000 & 1.000\\
0 & 0.0002 & 1.167 & -0.1667\\
0 & 0 & -5.251\times 10^4 & 7.504\times 10^3
\end{array}
\right].
\label{eq:no-pivot}
\end{equation}

Back substitution from \eqref{eq:no-pivot} gives
\[
z=\frac{7.504\times 10^3}{-5.251\times 10^4}=-0.1429,
\]
\[
0.0002y+1.167z=-0.1667
\qquad \Longrightarrow \qquad
y=0.5000,
\]
and then
\[
x-5y-z=1
\qquad \Longrightarrow \qquad
x=3.357.
\]
So Gaussian elimination with $4$-digit rounding and no pivoting gives
\[
\boxed{
\begin{pmatrix}
x\\
y\\
z
\end{pmatrix}
\approx
\begin{pmatrix}
3.357\\
0.5000\\
-0.1429
\end{pmatrix}}.
\]

\item Now we redo the same problem with partial pivoting. At each step, we choose the largest available entry in absolute value in the pivot column. In the first column, the largest entry in absolute value is $2.000$, so we swap rows $1$ and $3$ first:
\[
\left[
\begin{array}{ccc|c}
2.000 & -1.000 & 0 & 3.000\\
0.1667 & -0.8333 & 1.000 & 0\\
1.000 & -5.000 & -1.000 & 1.000
\end{array}
\right].
\]

\[
R_2\leftarrow R_2-0.08335R_1,\qquad R_3\leftarrow R_3-0.5000R_1,
\]
we get
\[
\left[
\begin{array}{ccc|c}
2.000 & -1.000 & 0 & 3.000\\
0 & -0.7500 & 1.000 & -0.2501\\
0 & -4.500 & -1.000 & -0.5000
\end{array}
\right].
\]

In the second column, the larger pivot is now $-4.500$, so we swap rows $2$ and $3$:
\[
\left[
\begin{array}{ccc|c}
2.000 & -1.000 & 0 & 3.000\\
0 & -4.500 & -1.000 & -0.5000\\
0 & -0.7500 & 1.000 & -0.2501
\end{array}
\right].
\]

\[
R_3\leftarrow R_3-0.1667R_2
\]
gives
\begin{equation}
\left[
\begin{array}{ccc|c}
2.000 & -1.000 & 0 & 3.000\\
0 & -4.500 & -1.000 & -0.5000\\
0 & 0 & 1.167 & -0.1668
\end{array}
\right].
\label{eq:with-pivot}
\end{equation}

Back substitution from \eqref{eq:with-pivot} gives
\[
z=\frac{-0.1668}{1.167}=-0.1429,
\]
\[
-4.500y-z=-0.5000
\qquad \Longrightarrow \qquad
y=0.1429,
\]
and then
\[
2.000x-y=3.000
\qquad \Longrightarrow \qquad
x=1.572.
\]
So with partial pivoting we get
\[
\boxed{
\begin{pmatrix}
x\\
y\\
z
\end{pmatrix}
\approx
\begin{pmatrix}
1.572\\
0.1429\\
-0.1429
\end{pmatrix}}.
\]

\item The exact answer from \eqref{eq:exact-solution} is
\[
\begin{pmatrix}
1.571428\ldots\\
0.142857\ldots\\
-0.142857\ldots
\end{pmatrix}.
\]
So the partial-pivoting answer from \eqref{eq:with-pivot} is very close, while the no-pivot answer from \eqref{eq:no-pivot} is pretty bad in the first two entries. The reason is the tiny pivot $0.0002$ that showed up in part (b). Dividing by such a small number created the enormous factor $4.500\times 10^4$, and that amplified the rounding errors. In part (c), partial pivoting avoids that by moving the largest available entry in the pivot column into the pivot position before eliminating, so the rounded arithmetic stays much more stable.
\end{enumerate}

\end{document}
